\documentclass[tikz, usenames, dvipsnames]{standalone}

\usepackage{tikz-3dplot}
\usepackage{xcolor}
\usepackage{pgfplots}
\usepackage{bm}

\usetikzlibrary{angles}
\usetikzlibrary{quotes}

\usepgfplotslibrary{fillbetween}

\renewcommand*{\vec}[1]{\bm{#1}}
\newcommand{\T}[2]{\vec{T}_{\vec{#1}}^{#2}}

\definecolor{color1}{RGB}{202, 0, 32}
\definecolor{color2}{RGB}{244, 165, 130}
\definecolor{color3}{RGB}{146, 197, 222}
\definecolor{color4}{RGB}{5, 113, 176}

\definecolor{arrowred}{RGB}{255, 0, 0}
\definecolor{arrowgreen}{RGB}{0, 170, 80}
\definecolor{arrowblue}{RGB}{0, 0, 255}

\tikzset{
    axis/.style={-latex, black, dashed}, 
	vector/.style={-latex, thick}, 
	vector guide/.style={-latex, thick},
	vector help one/.style={
		very thin,
		dash pattern={on 2pt off 1.5pt on 0.5pt off 1.5pt},
	},
	vector help two/.style={
		very thin,
		dash pattern={on 2pt off 2pt},
	},
	vector coord/.style={-latex},
}


\begin{document}

\tdplotsetmaincoords{54}{130}

\begin{tikzpicture}[
	scale=1,
	tdplot_main_coords,
]

\coordinate (O) at (0,0,0);

\def\x{3.5}
\def\sq{1.73}
\def\smallvecx{0.8}

\phantompathxy

\drawyz{
	(0, 3.5, 0)
	(0, 3.36, 0.24)
	(0, 3.14, 0.54)
	(0, 2.84, 0.90)
	(0, 2.50, 1.14)
	(0, 2.10, 1.26)
	(0, 1.64, 1.50)
	(0, 1.14, 2.06)
	(0, 0.62, 2.90)
	(0, 0.18, 3.36)
	(0, 0, 3.5)
};

\drawzx{
	(3.5, 0, 0)
	(3.08, 0, 0.70)
	(2.70, 0, 1.02)
	(2.38, 0, 1.10)
	(2.08, 0, 1.12)
	(1.76, 0, 1.22)
	(1.42, 0, 1.66)
	(1.02, 0, 3.25)
	(0.58, 0, 3.38)
	(0.16, 0, 3.45)
	(0, 0, 3.5)
};

% Same as E7 with (A) coordinates copied to (B) and x <-> z.

\def\Bx{0}
\def\By{0}
\def\Bz{2.3}

\coordinate (B) at (\Bx, \By, \Bz);

\path (O) + (\smallvecx, 0, 0) coordinate (O1);
\path (O) + (0, \smallvecx, 0) coordinate (O2);
\path (O) + (0, 0, \smallvecx) coordinate (O3);

\path (B) + (\smallvecx, 0, 0) coordinate (B1);
\path (B) + (0, -\smallvecx, 0) coordinate (B2);
\path (B) + (0, 0, \smallvecx) coordinate (B3);

\drawaxes

\draw[
	vector, ProcessBlue
] (O) -- (B) node [
	anchor=west,
	xshift=0pt, yshift=-4pt,
] {\textcolor{black}{\(\T{A}{}\)}};

\drawtriad{O}
\drawtriad{B}

\end{tikzpicture}
\end{document}
